\chapter{Zakończenie}

\section{Podsumowanie celu pracy}

Celem pracy było zbudowanie i ocena modelu probabilistycznego dla profesjonalnych meczów
\textit{League of Legends}. Model miał wykorzystywać informacje dostępne przed meczem, w
szczególności systemy rankingowe zawodników i drużyn, historyczny kontekst sportowy oraz format
serii. Drugim celem była ocena, czy taki model może osiągać jakość predykcyjną konkurencyjną
względem prawdopodobieństw implikowanych przez kursy bukmacherskie.

Zakres pracy był świadomie ograniczony do predykcji probabilistycznej. Kursy bukmacherskie pełniły
funkcję benchmarku rynkowego, ale nie analizowano strategii obstawiania, zwrotu finansowego ani
zarządzania kapitałem. Dzięki temu głównym kryterium oceny pozostała jakość prawdopodobieństw,
przede wszystkim LogLoss.

\section{Najważniejsze wyniki}

Pierwszym istotnym wynikiem jest przewaga rankingów zawodniczych nad rankingami drużynowymi. W
kontekście esportu oznacza to, że aktualny skład lepiej opisuje siłę zespołu niż sama historyczna
nazwa organizacji. Najlepszym pojedynczym systemem rankingowym okazał się Player Glicko-2.

Drugim wynikiem jest skuteczność metamodelu W20-Binomial. Połączenie predykcji rankingowych, miar
niepewności, historycznych cech kroczących oraz korekty formatu serii poprawiło jakość
probabilistyczną względem pojedynczego rankingu. Najlepszym estymatorem dla tego zestawu cech była
regularyzowana regresja logistyczna ElasticNet.

Trzecim wynikiem jest przewaga finalnego wariantu Sym-Cal LR-ElasticNet-W20-Binomial. Końcowa
symetryzacja kolejności drużyn i kalibracja Platta poprawiły surowy wariant regresji logistycznej w
sposób niewielki, lecz stabilny. Na wspólnej próbie sportowo-rynkowej model finalny uzyskał
najniższy LogLoss spośród analizowanych wariantów i benchmarków.

\section{Weryfikacja hipotez}

Hipoteza pierwsza, zgodnie z którą rankingi zawodnicze osiągają lepszą jakość predykcji niż rankingi
drużynowe, została potwierdzona. Wynik ten jest zgodny ze specyfiką \textit{League of Legends},
gdzie transfery i zmiany składów istotnie wpływają na bieżącą siłę drużyny.

Hipoteza druga, dotycząca przewagi metamodelu nad pojedynczym rankingiem, również została
potwierdzona. Player Glicko-2 był silnym punktem odniesienia, ale metamodel rankingowo-kontekstowy
osiągnął lepsze wartości LogLoss i Brier Score.

Hipoteza trzecia, zakładająca przydatność historycznego kontekstu W20, została potwierdzona w sposób
umiarkowany. Cechy kroczące nie zastąpiły ratingów, lecz uzupełniły je o informacje o formie, tempie
gry, ekonomii i stylu drużyny.

Hipoteza czwarta, odnosząca się do konkurencyjności względem rynku bukmacherskiego, została
potwierdzona na analizowanej próbie. Finalny model uzyskał niższy LogLoss niż Market Open i Market
Close, a miesięczny bootstrap blokowy wskazał dodatnie przedziały ufności. Wniosek ten dotyczy
jakości predykcyjnej, nie wyniku finansowego.

\section{Odpowiedź na pytanie badawcze}

Odpowiedź na główne pytanie badawcze jest pozytywna. Model oparty na rankingach zawodników,
historycznym kontekście drużyn, korekcie formatu serii, symetryzacji predykcji i kalibracji Platta
może skutecznie przewidywać wyniki profesjonalnych meczów \textit{League of Legends}. Na wspólnej
próbie z rynkiem bukmacherskim skalibrowany symetryczny LR-ElasticNet-W20-Binomial osiągnął niższy
LogLoss niż benchmarki Market Open i Market Close.

Jednocześnie odpowiedź ta ma ograniczony zakres. Wyniki odnoszą się do danych historycznych,
zastosowanej procedury walidacji i meczów posiadających kursy bukmacherskie. Nie stanowią dowodu
opłacalności strategii zakładów ani gwarancji utrzymania przewagi w przyszłych sezonach.
Najważniejszym rezultatem pracy jest więc pokazanie, że starannie skonstruowany model
rankingowo-kontekstowy może stanowić silny i interpretowalny model probabilistyczny dla
profesjonalnego esportu.
